\section{Introduction}
\label{sec_introduction}

\begin{figure*}[t]{
    \centering
        \begin{subfigure}{0.4\textwidth}
            \includegraphics[width=\textwidth]{Figures/Method/algorithm.pdf}
            \caption{Schematic overview}
        \end{subfigure}    
        \begin{subfigure}{0.45\textwidth}
            \includegraphics[width=\textwidth]{Figures/Method/moons_trial_0_for_publication_v2.pdf}
            \caption{2-dimensional trajectory}
        \end{subfigure}
    \caption{An overview of the safe denoiser. (a) The safe denoiser $\mathbb{E}_{\text{safe}}$ negates the direction of the unsafe denoiser $\mathbb{E}_{\text{unsafe}}$ from the data denoiser $\mathbb{E}_{\text{data}}$. 
    (b) Trajectories from data denoiser and safe denoiser, starting from the same initial point far from the data distribution, reveal distinct paths: while the sample path from the data denoiser falls into the unsafe region, the trajectory from the safe denoiser successfully avoids it.}
    \label{fig:schematic}
}
\end{figure*}

Diffusion models (DMs) have become leading generative models, excelling in generation tasks like text-to-image~\cite{rombach2022high}, audio~\cite{kong2021diffwave}, video~\cite{bartal2024lumiere}, and protein design~\cite{watson2023denovo}, thanks to their flexible and controllable sampling~\cite{dhariwal2021diffusion, ho2021classifier}. However, growing concerns over unsafe content — such as not-safe-for-work(NSFW) imagery (\figref{thumbnail_1}), copyright violations, and potential misuse — highlight the need for safety. The key challenge is mitigating these risks without compromising model utility or creativity.

Mainstream mitigation strategies for issues like NSFW content or unwanted concept removal rely on text-based guidance~\cite{gandikota2023erasing, ban2024understanding, yoon2024safree} or fine-tuning for unlearning~\cite{gandikota2023erasing, gandikota2023unified, gong2024reliable, kim2024race}. Text-based methods require iterative, expert-crafted negative prompts~\citep{schramowski2023safe}, which may not generalize well, while fine-tuning is resource-intensive and risks catastrophic forgetting or degraded performance on desired tasks.

Other significant safety concerns involve the DMs' capacity to reproduce copyrighted content and their generation of data pertaining to individuals in \figref{thumbnail_2} who wish to be excluded. These issues are often linked to the models' remarkable ability to memorize training data \cite{carlini2023extracting}. While techniques like differentially private training \cite{dockhorn2023differentially, liu2024differentially} can formally limit memorization by adding noise during the training process, they often result in a noticeable degradation in generation quality, which can be particularly prohibitive for applications demanding high-fidelity outputs.

We propose a \textit{safe denoiser} (defined in Definition~\ref{thm:def}) that modifies sampling trajectories such that the resulting samples are drawn from a safe distribution 
(shown in \figref{schematic}). The intuion comes from our \thmref{safe}, 
% which relates safe and unsafe denoised samples, 
where the safe denoiser steers generation away from unsafe regions, ensuring theoretical safety. We develop a practical algorithm (\algoref{safer}) based off of our theorem, which can be used standalone or combined with negative prompting to enhance safety in text-to-image generation. 
% Empirically, 
Our method achieves state-of-the-art performance on concept erasing, class removal, and unconditional image generation tasks.